\documentclass[11pt,a4paper]{article}
\usepackage[T1]{fontenc}
\usepackage[utf8]{inputenc}
\usepackage[brazilian]{babel}
\usepackage{lmodern}
\usepackage[a4paper,margin=2.5cm]{geometry}
% espaçamento 1,5 (mesmo fator que o \onehalfspacing do pacote setspace usa em 11pt)
\pagestyle{empty}
\linespread{1.213}\selectfont
\setlength{\parindent}{0pt}

\begin{document}

\begin{center}
{\large\bfseries Geração de bateria expressiva condicionada por estilo com Transformer e \textit{classifier-free guidance}}

\vspace{0.6em}
Renato Gabino Diniz, Isadora Cunha Pimentel, Vinicius Pinho Galvão

Universidade Federal de Minas Gerais (UFMG) -- Departamento de Ciência da Computação\\
DCC831 -- IA Generativa para Música -- 2026/2
\end{center}

\vspace{0.4em}
\textbf{Resumo.}
Em bateria, o que distingue um estilo do outro está em grande parte no \textit{groove}: o microtiming e a dinâmica com
que cada nota é tocada, que se perdem quando a música é quantizada numa grade rítmica. Este trabalho aborda a geração
simbólica de bateria expressiva condicionada por estilo musical (rock, funk, latin, jazz, hip-hop e soul) a partir do
\textit{Groove MIDI Dataset}, cujas performances em 4/4 somam 491 gravações, 12,7 h e cerca de 21 mil compassos,
com forte desbalanceamento entre estilos. As performances são representadas com a tokenização PerTok, que combina uma
grade de semicolcheia com um \textit{token} de desvio temporal por nota e 16 níveis de \textit{velocity}, preservando
o \textit{swing} com erro de \textit{onset} de cerca de 2 ms. O modelo principal é um Transformer
\textit{decoder-only} implementado em PyTorch, no qual o estilo entra como um \textit{embedding} somado a cada posição
e é substituído por um estilo nulo em 20\% das sequências de treino, o que permite aplicar \textit{classifier-free
guidance} na geração para regular a intensidade do estilo e trocá-lo no meio da música; a amostragem combina
temperatura, \textit{top-k} e \textit{top-p}, e a amostragem ponderada por estilo mitiga o desbalanceamento dos dados.
Como \textit{baseline}, treinamos uma cadeia de Markov de ordem $n$ por estilo, com \textit{backoff}, sobre os mesmos
\textit{tokens}, com a ordem escolhida pela verossimilhança na validação. Os dois sistemas geram 100 amostras de
8 compassos por estilo, comparadas com janelas humanas do mesmo estilo por métricas de \textit{groove} --
taxa de tempos vazios, consistência de \textit{groove}, consistência de padrão, densidade, \textit{swing} e desvio de
microtiming -- confrontando distribuições, e não apenas médias, por área de sobreposição e divergência KL. A eficácia
do condicionamento é verificada por uma matriz que mede a distância entre cada estilo gerado e cada estilo real, sem
recorrer a um classificador, e o papel do microtiming é examinado comparando cada amostra com sua versão quantizada.
O trabalho resulta em um \textit{pipeline} reprodutível de ponta a ponta (tokenização, treino, geração, renderização em
áudio e avaliação) e discute os limites do modelo na estrutura de longo prazo e nos estilos com poucos dados.

\end{document}
